% Generated draft: requires mathematical and source review.
\documentclass[11pt]{article}
\usepackage[margin=25mm]{geometry}
\usepackage{fontspec}
\setmainfont{Latin Modern Roman}
\usepackage{amsmath,amssymb,mathtools,unicode-math}
\setmathfont{Latin Modern Math}
\usepackage{enumitem,needspace,xurl,hyperref}
\setlength{\parindent}{0pt}
\setlength{\parskip}{5pt}
\setlength{\emergencystretch}{3em}
\newcommand{\N}{\mathbb{N}}
\newcommand{\Z}{\mathbb{Z}}
\newcommand{\Q}{\mathbb{Q}}
\newcommand{\R}{\mathbb{R}}
\newcommand{\C}{\mathbb{C}}
\newcommand{\F}{\mathbb{F}}
\newcommand{\A}{\mathbb{A}}
\newcommand{\PP}{\mathbb{P}}
\newcommand{\E}{\mathbb{E}}
\newcommand{\eps}{\varepsilon}
\newcommand{\dd}{\,\mathrm{d}}
\begin{document}
\section*{UnsolvedMath ID: 1149}
\textbf{Source problem code:} catalogue entry
\subsection{Definitions and mathematical statement}
Let $C=[0,1]^3\subset\R^3$ be the unit solid cube. A smaller positively homothetic copy of $C$ is a set
\[
 C_i=t_i+\lambda_i C=\{t_i+\lambda_i x:x\in C\},
 \qquad t_i\in\R^3,\quad 0<\lambda_i<1.
\]
The interior is taken in $\R^3$. The literal covering question is whether there exist an integer $m\geq1$, translation vectors $t_1,\ldots,t_m$, and ratios $\lambda_1,\ldots,\lambda_m\in(0,1)$ such that
\[
 C\subseteq\bigcup_{i=1}^{m}C_i,
 \qquad
 \operatorname{int}(C_i)\cap\operatorname{int}(C_j)=\varnothing
 \quad(1\leq i<j\leq m).
\]
The word covering only requires the displayed inclusion. If the smaller cubes are also required to lie inside $C$, the covering condition is equivalently $C=\bigcup_i C_i$. Boundary intersections are allowed: the supplied statement requires disjoint interiors, not disjoint closed cubes.

\textbf{Source qualification.} The literature triage dated 2026-08-17 reports that this literal question has an elementary affirmative answer, including the stronger inside-$C$ interpretation: coordinate subdivision gives $2^3=8$ congruent cubes of homothety ratio $1/2$ with disjoint interiors covering $C$. It reports the analogous $2^d$ subdivision for a $d$-dimensional cube. This is a supplied status qualification, not a new open conjecture or a proof offered here.

The imported title calls the question ``The Knaster Problem,'' and the background says that a two-dimensional version was proved impossible by Mycielski in 1979 while the three-dimensional version remains open. Those assertions are incompatible with the literal disjoint-interior covering formulation and the elementary subdivision reported in the same record. The background also suggests links to measure theory, geometric coverings, and Banach--Tarski-like phenomena, but supplies no authoritative reference establishing this identification. The triage explicitly flags a likely formulation defect and says any different intended Knaster problem must be recovered. No omitted condition is guessed here. The registry label \emph{open} and dated triage label \emph{solved} are therefore recorded as conflicting source classifications, with the latter referring only to the literal question.
\subsection{Short English statement}
Ask whether a solid cube is covered by finitely many smaller homothetic cubes whose interiors do not overlap. The supplied dated review says ordinary subdivision answers the literal question affirmatively. Its identification with an open Knaster problem and the claimed two-dimensional impossibility are therefore flagged as source defects requiring recovery of the intended statement.
\subsection{Sources}
\begin{itemize}
\item UnsolvedMath local registry, data/unsolvedmath/problems.json, permanent ID 1149, source code catalogue entry, title ``The Knaster Problem.'' The statement, contradictory historical background, and literature triage dated 2026-08-17 are the supplied evidence.
\item Supplied review reference: ``Elementary cube subdivision argument (checked 2026-08-17),'' linked to \url{https://en.wikipedia.org/wiki/Cube}. The review identifies standard half-side coordinate subdivision as evidence for the literal covering question. This general cube reference does not identify an authoritative historical Knaster formulation or substantiate the Mycielski attribution; no such citation is supplied.
\item Ulam AI and the UnsolvedMath Contributors, \emph{UnsolvedMath: A Curated Collection of Open Mathematics Problems}. Dataset curation and original metadata are supplied under CC BY 4.0; referenced source material retains its own terms. No external source was independently consulted for this draft.
\end{itemize}
\end{document}
